\begin{helper}
Let $n\ge1$ be an integer, let $1<p<\infty$, let
$A\colon\mathbb{Z}\to\mathbb{C}$, let $m\ge0$ be an integer and let
$\xi\in\mathbb{C}$ with $\xi\ne0$.  Write $k=n+m$, $A_k=A(k)$,
$\lambda_k=\lambda_k(n,p)$ as in \noderef{LambdaCoef},
$\epsilon_k=\epsilon_k(n,p)$ as in \noderef{EpsCoef}, and let
$T_k=T_k(n,p,A,\cdot)$ be the $k$-th summand of the series (19) as in
\noderef{Thm2Term}.  Then
\[
  \Bigl|\frac{\partial T_k}{\partial\xi}(\xi)\Bigr|
  \le\bigl(4+\sqrt{p-1}\bigr)(n+m)\,|A_{n+m}|\,|\xi|^{\lambda_{n+m}-1}
\]
and
\[
  \Bigl|\frac{\partial T_k}{\partial\bar\xi}(\xi)\Bigr|
  \le\bigl(4+\sqrt{p-1}\bigr)(n+m)\,|A_{n+m}|\,|\xi|^{\lambda_{n+m}-1},
\]
with $\partial/\partial\xi$ as in \noderef{DZ} and $\partial/\partial\bar\xi$
as in \noderef{DZbar}.

The content of the statement is that the constant $4+\sqrt{p-1}$ does not
depend on $m$: the only $m$-dependence left on the right is the explicit
factor $n+m$ and the explicit power $|\xi|^{\lambda_{n+m}-1}$.  That
uniformity comes entirely from the Remark's estimates (23)--(25), carried by
\noderef{LambdaEpsEstimates}: $0\le\lambda_{n+m}\le(n+m)\sqrt{p-1}$ and
$|\epsilon_{n+m}|\le1$.  Without the \emph{upper} bound on $\lambda_k$ the
coefficients of the closed form would grow in $k$ at an unknown rate and no
uniform constant would exist.

This is the pointwise estimate the $W^{1,2}$ convergence argument needs before
it can sum over $k$: on a disc of radius $\rho<1$ the factor
$|\xi|^{\lambda_{n+m}-1}$ decays geometrically in $m$, so this bound makes the
tail of the differentiated series dominated by a convergent one, which is what
allows (19) to be differentiated term by term in the $W^{1,2}$ norm.

The family is indexed by $m$ rather than by $k\ge n$ because that is the shape
in which \noderef{LambdaEpsEstimates} is used throughout: all three of its
estimates are read as statements about $\lambda_{n+m}$ and $\epsilon_{n+m}$.
The constant $4$ is not sharp; the route below gives $3$.
\end{helper}
\begin{proof}
Throughout, $k=n+m$, $r=|\xi|>0$, and
\[
  s=\lambda_k+n-k=\lambda_{n+m}-m,\qquad
  F=A_k\xi^{k}+\epsilon_k\overline{A_k}\,\bar\xi^{\,k},
\]
so that $T_k(\xi)=F\,|\xi|^{s}\xi^{-n}$ by \noderef{Thm2Term}.

\emph{Step 0: the three estimates used.}  From
\noderef{LambdaEpsEstimates} we record, for every $m\ge0$,
\begin{equation}\tag{1}
  0\le\lambda_{n+m},\qquad
  \lambda_{n+m}\le\bigl((n:\mathbb{R})+m\bigr)\sqrt{p-1},\qquad
  |\epsilon_{n+m}|\le1 .
\end{equation}
For $m=0$ these are (23), $\lambda_n=0$ and $\epsilon_n=0$, together with
$0\le n\sqrt{p-1}$; for $m\ge1$ we have $n+1\le n+m$, so the first and second
are (25) and the third is (24).

\emph{Step 1: the closed form.}  By \noderef{Thm2TermWirtingerDerivs}, applied
at the index $k=n+m$ and the point $\xi\ne0$,
\begin{equation}\tag{2}
  \frac{\partial T_k}{\partial\xi}(\xi)
  =\frac{|\xi|^{s}\xi^{-n}}{\xi}
    \Bigl(k\,A_k\xi^{k}+\Bigl(\frac{s}{2}-n\Bigr)F\Bigr),
  \qquad
  \frac{\partial T_k}{\partial\bar\xi}(\xi)
  =\frac{|\xi|^{s}\xi^{-n}}{\bar\xi}
    \Bigl(k\,\epsilon_k\overline{A_k}\,\bar\xi^{\,k}+\frac{s}{2}F\Bigr).
\end{equation}
That node requires no constraint on $n$, $p$ or $k$, only $\xi\ne0$.

\emph{Step 2: the modulus of the common prefactor.}  Since $r>0$, real and
integer powers of $r$ multiply by adding exponents, and
$\bigl||\xi|^{s}\bigr|=r^{s}$, $|\xi^{-n}|=r^{-n}$, $|\xi|=|\bar\xi|=r$.  Hence
\begin{equation}\tag{3}
  \Bigl|\frac{|\xi|^{s}\xi^{-n}}{\xi}\Bigr|
  =\Bigl|\frac{|\xi|^{s}\xi^{-n}}{\bar\xi}\Bigr|
  =r^{s-n-1}=r^{\lambda_k-k-1},
\end{equation}
using $s-n-1=(\lambda_k+n-k)-n-1=\lambda_k-k-1$.

\emph{Step 3: the two brackets.}  Since $|\overline{A_k}|=|A_k|$,
$|\bar\xi|=r$ and $\epsilon_k$ is real, the two terms of $F$ have moduli
$|A_k|r^{k}$ and $|\epsilon_k||A_k|r^{k}$, so by the triangle inequality and
(1),
\begin{equation}\tag{4}
  |F|\le\bigl(1+|\epsilon_k|\bigr)|A_k|\,r^{k}\le2|A_k|\,r^{k}.
\end{equation}
Therefore, again by the triangle inequality and (4),
\[
  \Bigl|k\,A_k\xi^{k}+\Bigl(\frac{s}{2}-n\Bigr)F\Bigr|
  \le k\,|A_k|r^{k}+\Bigl|\frac{s}{2}-n\Bigr|\cdot2|A_k|r^{k}
  =\bigl(k+|s-2n|\bigr)|A_k|\,r^{k},
\]
\[
  \Bigl|k\,\epsilon_k\overline{A_k}\bar\xi^{\,k}+\frac{s}{2}F\Bigr|
  \le k\,|\epsilon_k||A_k|r^{k}+\Bigl|\frac{s}{2}\Bigr|\cdot2|A_k|r^{k}
  \le\bigl(k+|s|\bigr)|A_k|\,r^{k},
\]
where $2\bigl|\tfrac{s}{2}-n\bigr|=|s-2n|$, $2\bigl|\tfrac{s}{2}\bigr|=|s|$,
and $|\epsilon_k|\le1$ was used in the second display.

\emph{Step 4: the two coefficient sizes are at most
$(4+\sqrt{p-1})(n+m)$.}  Write $N=(n:\mathbb{R})+m$, so $k=N$ as a real number
and, by (1), $0\le\lambda_k\le N\sqrt{p-1}$.  Since $\lambda_k\ge0$ and
$m\ge0$, the number $s=\lambda_k-m$ satisfies
$-(\lambda_k+m)\le s\le\lambda_k+m$, i.e.
\begin{equation}\tag{5}
  |s|\le\lambda_k+m\le N\sqrt{p-1}+m .
\end{equation}
Hence, using $m\le N$ and $n\le N$,
\[
  k+|s|\le N+N\sqrt{p-1}+m\le 2N+N\sqrt{p-1}
  =N\bigl(2+\sqrt{p-1}\bigr)\le N\bigl(4+\sqrt{p-1}\bigr),
\]
and, since $|s-2n|\le|s|+2n\le N\sqrt{p-1}+m+2n\le N\sqrt{p-1}+2N$,
\[
  k+|s-2n|\le N+N\sqrt{p-1}+2N=N\bigl(3+\sqrt{p-1}\bigr)
  \le N\bigl(4+\sqrt{p-1}\bigr).
\]
(The second chain is where the constant is largest, and it is the reason $3$
would already do.)

\emph{Step 5: collecting.}  Multiply the bound of Step 3 by the modulus (3) of
the prefactor, which is legitimate because both factors are nonnegative, and
use $r^{k}\,r^{\lambda_k-k-1}=r^{\lambda_k-1}$ (again $r>0$).  For the
$\partial/\partial\xi$ derivative this gives, by (2),
\[
  \Bigl|\frac{\partial T_k}{\partial\xi}(\xi)\Bigr|
  \le\bigl(k+|s-2n|\bigr)|A_k|\,r^{\lambda_k-1}
  \le\bigl(4+\sqrt{p-1}\bigr)N\,|A_k|\,r^{\lambda_k-1},
\]
the last step by Step 4 together with $|A_k|r^{\lambda_k-1}\ge0$; and for the
$\partial/\partial\bar\xi$ derivative, identically,
\[
  \Bigl|\frac{\partial T_k}{\partial\bar\xi}(\xi)\Bigr|
  \le\bigl(k+|s|\bigr)|A_k|\,r^{\lambda_k-1}
  \le\bigl(4+\sqrt{p-1}\bigr)N\,|A_k|\,r^{\lambda_k-1}.
\]
Since $N=(n:\mathbb{R})+m$, $k=n+m$ and $A_k=A_{n+m}$, these are the two
asserted inequalities.
\end{proof}
